{\Large \noindent Abstract}

Additive Manufacturing (AM) is a relatively new method of manufacturing in which material is added layer by layer.  Initially this happened for resins and plastics, with Vat Photopolymerization (VPP) being one of the first methods.  This method involves resins fusing as a result of being exposed to light or lasers.  Fused Deposition Modeling (FDM) stacks extruded plastics on top of each other to create parts.  More recently, metals have been added to the library of materials that can be made into parts via AM.
\newline \indent Powder bed fusion is a method of three-dimensional (3D) printing metals that has risen to prominence recently since its cost has come down.  Companies like Czinger Automotive have combined topology optimization and metal 3D printing to create performance parts that are stronger and lighter than what is possible through traditional manufacturing methods.  It has also found its way into aerospace; making turbine parts out of superalloys to be sent into the harshest environments imaginable. \newline \indent The most common form of metal 3D printing is Laser Powder Bed Fusion (LPBF).  This is a process in which powdered metal is fused together layer by layer with a laser to produce a monolithic part.  However, the fusion process leaves behind residual stresses in the finished part due to the uneven distribution of heat in the powder and melt pool.
\newline \indent The goal of this project is to create a tool that can minimize the residual stress in a part by finding the ideal orientation in 3D space to print the part in question.  To accomplish this, Solidworks will be used to rotate parts, and ANSYS will be used to find the residual stress in the part in its particular orientation.  To control these two programs, the tool will be created in Python.  This tool will be able to rotate the part in Solidworks, find the residual stress in ANSYS, and intelligently iterate to find the orientation with the least residual stress. 